\chapter{Точный Casimir--Hodge профиль и циклический тяжёлый гессиан}
\label{ch:v7-exact-profile-hodge-cycle-unification-gate}

\section{Возврат отброшенных Casimir-весов}

Предыдущий гейт применял точный Gaussian только к физическому оператору
$\Phi$. Формульная археология показала, что Hodge-вакуум живёт на другом
производном объекте --- отображении момента. Для заряженного ребра с
собственным значением полного gauge-Casimir $c_e>0$ точный Hodge-тепловой
вклад имеет нулевую кривизну
\begin{equation}
 h_e(t)=8tc_e e^{-tc_e^2},\qquad t>0.
 \label{eq:v7-exact-unification-hodge-mass}
\end{equation}
Здесь общий масштаб положен единичным после прежней массовой калибровки.

Выжившие ветви различаются уже выведенными значениями
\begin{equation}
 c_d=\frac85,\qquad c_W=\frac9{10}.
 \label{eq:v7-exact-unification-casimir-values}
\end{equation}
Это не новый порог: оба числа являются собственными значениями одного
оператора $\mathcal C_G$ на полных калибровочных мультиплетах.

\section{Один формальный градуированный профиль}

Минимальный кандидат записывается как один градуированный след на прямой
сумме рёберного Hodge- и физического носителей:
\begin{equation}
 \mathcal S_t^{\rm gr}
 =-\Tr_{\rm edge}e^{-t\mathfrak m_C^2}
  +\Tr_{\rm phys}e^{-t\Phi^2}.
 \label{eq:v7-exact-unification-graded-profile}
\end{equation}
Обе части используют одну и ту же тепловую шкалу и единичные следовые
кратности. Пока это формальный общий профиль, а не доказанная единая
физическая суперсвязность.

На двадцатимерном вещественном тяжёлом пространстве его гессиан равен
\begin{equation}
 H_{\rm tot}(t)=H_{\rm phys}(t)
 +\operatorname{diag}\bigl(h_e(t)\bigr).
 \label{eq:v7-exact-unification-total-hessian}
\end{equation}

\section{Открытое положительное окно}

Точный спектральный расчёт даёт при $t=1$
\begin{equation}
 \lambda_{\min}H_{\rm tot}(1)=1.03081235398>0,
 \qquad (n_-,n_0,n_+)=(0,0,20).
 \label{eq:v7-exact-unification-benchmark}
\end{equation}
Таким образом, ненулевые Casimir-веса действительно снимают отрицательную
weak-моду, не разрушая down-блок.

Логарифмический проход по $601$ точке и последующее уточнение корня дают
\begin{equation}
 0<t<t_*\quad\Longrightarrow\quad H_{\rm tot}(t)>0,
 \qquad t_*\simeq2.36617,
 \label{eq:v7-exact-unification-positive-window}
\end{equation}
где точное численное значение $t_*$ хранится в машинном сертификате.

\section{Оставшийся корневой барьер}

Положительный тяжёлый гессиан ещё не делает фон вакуумом. Физический
Gaussian сохраняет ненулевые производные по семи корневым амплитудам;
Hodge-тепловой вклад в своём минимуме имеет нулевой градиент и сам по себе
их не сокращает. Кроме того, прямосуммовая запись
\eqref{eq:v7-exact-unification-graded-profile} ещё не выведена квадратом
одной Real-суперсвязности.

\begin{theorem}[Casimir-взвешенный тяжёлый проход]
При использовании полных собственных значений gauge-Casimir, а не только
проектора на его ядро, точный Hodge-тепловой вклад стабилизирует совместный
down--weak тяжёлый сектор в непустом открытом интервале общей тепловой
шкалы. В частности, при $t=1$ все двадцать собственных значений положительны.
\end{theorem}

\begin{nogocriterion}[Тяжёлый проход ещё не равен общему родителю]
Нельзя объявлять формальную разность двух следов единым действием, пока она
не получена из одной градуировки, одной Real-суперсвязности и одного
физического следа. Нельзя также игнорировать ненулевые корневые производные.
\end{nogocriterion}

\begin{protocol}
\begin{enumerate}
 \item полный Casimir слабой пары: \textbf{$9/10$};
 \item полный Casimir down-пары: \textbf{$8/5$};
 \item независимый Hodge-вес: \textbf{не введён};
 \item положительное тяжёлое окно: \textbf{есть};
 \item сигнатура при $t=1$: \textbf{$(0,0,20)$};
 \item стационарность корневого фона: \textbf{нет};
 \item единая физическая суперсвязность: \textbf{не выведена};
 \item итог: \textbf{существенный частичный проход};
 \item следующий гейт: \textbf{общий носитель и корневая стационарность}.
\end{enumerate}
\end{protocol}

Машинный сертификат сохранён в
\path{s2t_v7_exact_profile_hodge_cycle_unification_gate_results.json}.